\chapter{Proposed approach}\label{chap:Proposed approach}

To ensure the feasibility of applying Deep Reinforcement Learning (DRL), the dynamic retrieval and relocation scheduling problem is formulated as a Markov Decision Process (MDP), represented by a five-tuple$\{\mathcal{S}, \mathcal{A}, R, \gamma, p\}$where$\mathcal{S}$ denotes the observation (state) space,$\mathcal{A}$ denotes the action space, $R$ denotes the reward function, $\gamma$ is the discount factor used to balance immediate and long-term rewards, and $p$ denotes the transition probability that characterizes the dynamics of the environment.

A fundamental assumption underlying an MDP is the Markov property, which states that the future evolution of the system depends only on the current state and is independent of the preceding state-action history, given the current state. On the one hand, this property eliminates the need for the agent to retain the complete historical trajectory from the initial state to the current state. On the other hand, provided that the current state $s$ contains sufficient information, the agent can determine an action $a$ solely based on $s$, thereby making large-scale computation and policy optimization tractable.

Within the environment formulated as an MDP, the DRL agent continuously interacts with the environment by selecting actions and receives feedback through the reward function to optimize its policy. Specifically, at each decision point $t$, the agent observes the current state $s_t$ from the environment and uses $s_t$ as the input to its policy $\pi$, which selects an action $a_t$. The execution of $a_t$ causes the environment to transition from $s_t$ to the subsequent state $s_{t+1}$, while the environment provides the agent with a reward $r_t$. Through repeated interactions with the environment, the agent progressively improves its policy, with the objective of maximizing the expected cumulative reward over an episode.

In this study, the agent is constructed using a hybrid neural network architecture that integrates a Graph Neural Network (GNN) and a Transformer, and the network parameters are optimized using the Proximal Policy Optimization (PPO) algorithm. The graph-based state representation and the scalable action space are introduced in Sections~4.1 and~4.2, respectively. The specific neural network architecture is described in Section~4.3. Section~4.4 provides a detailed description of the reward function design. The PPO mechanism for updating the network parameters is presented in Section~4.5.

\section{State representation}
At each upper-level decision point, the warehouse is represented as an undirected graph and provided as input to the deep reinforcement learning (DRL) agent. Each node in the graph corresponds to a physical position in the warehouse layout, including a storage location, a cross-aisle position, or an I/O point. Edges between adjacent nodes represent the physical connectivity of the warehouse. Adjacent storage locations within the same storage lane are sequentially connected, the entrance of each lane is connected to the cross-aisle, and cross-aisle nodes are connected according to the physical warehouse layout. In this way, the graph explicitly captures the spatial topology of the multi-deep warehouse as well as the structural dependencies induced by storage depth within each lane. It should be noted that items are not modeled as separate graph nodes. Instead, each storage location is represented by a fixed node, while whether the location currently contains a target item, a non-target item, or is empty is encoded by its node features.

At each state, every node is described by a nine-dimensional feature vector. The first three dimensions jointly represent the physical node type. The first dimension indicates whether the node is a storage location, the second dimension indicates whether it is a cross-aisle node, and the third dimension indicates whether it is an I/O point. The fourth and fifth dimensions represent the vertical and horizontal coordinates of the node in the warehouse layout, respectively, thereby providing the agent with spatial position information.

The sixth dimension describes the current task attribute of a storage location. If the node contains an activated target item, this dimension represents the due date of the corresponding item. If the node contains a non-target item that has not yet been activated, the value is set to $-1$. If the storage location is empty, or if the node is not a storage location, the value is set to $0$. Therefore, this feature jointly provides information about item occupancy, whether the stored item is currently a target item, and the due-date information of active target items.

The seventh dimension represents the current position and carrying state of the shuttle. If the shuttle is located at the node while carrying an item, the value is set to $1$. If the shuttle is located at the node without carrying an item, the value is set to $-1$. All other nodes take the value $0$. This feature allows the agent to identify the shuttle position from the current state and to distinguish whether the current upper-level decision is made during the unloaded item-selection stage or during the loaded delivery or relocation stage.

The eighth dimension represents the current environment time. Together with the due dates of the target items, this information characterizes the temporal status of the current scheduling problem and enables the agent to distinguish the urgency of different target items at the current decision point.

The ninth dimension provides additional scheduling information associated with the adopted reward model, and its specific meaning therefore varies across reward formulations. Under the original tardiness reward, this dimension represents the buffered tardiness snapshot of an active target item. This quantity is constructed from the due date of the target item, its static shortest path to the nearest I/O point, and the current time, thereby reflecting its tardiness status while accounting for the potential transportation time. For non-target items, empty storage locations, cross-aisle nodes, and I/O points, this feature is set to $0$.

For the three lower-bound-based reward models, the ninth dimension no longer represents buffered tardiness and is instead replaced by target-specific lower-bound information derived from the current physical warehouse state. Under \texttt{simple\_shuttle\_LB}, the ninth dimension represents the required-work lower bound $q$ of the corresponding target item. This quantity consists of the retrieval-operation lower bound required to complete the target item together with the relocation-operation lower bounds of the blocking non-target items currently associated with that target, and therefore estimates the minimum amount of work required to complete the target item.

Under \texttt{mixed\_shuttle\_LB} and \texttt{topology\_shuttle\_LB}, the ninth dimension represents the lane-prefix cumulative work lower bound $R$ of the corresponding target item. Within each storage lane, $R$ is obtained by cumulatively aggregating the $q$ values from the first target item in the lane to the current target item. Consequently, $R$ further encodes the lane-level precedence relationships induced by the multi-deep storage topology. For all lower-bound-based reward models, the ninth dimension contains the corresponding lower-bound information only for currently active target-item nodes and is set to $0$ for all other nodes.

Accordingly, although all reward models retain the same nine-dimensional node-feature structure, the ninth dimension follows different node-feature schemas according to the adopted reward model. The tardiness reward uses buffered tardiness, \texttt{simple\_shuttle\_LB} uses $q$, whereas \texttt{mixed\_shuttle\_LB} and \texttt{topology} \texttt{\_shuttle\_LB} use $R$. This design keeps the dimensionality of the state input unchanged while allowing the scheduling lower-bound information associated with the corresponding reward formulation to be incorporated directly into the agent's state representation.

In addition to the node-level features, the current state contains an independent graph-level dynamic context that indicates whether the dynamic target activation process is currently enabled. This variable takes the value $1$ when the dynamic activation process has not exceeded its predefined termination time and at least two non-target items remain unactivated in the warehouse; otherwise, it takes the value $0$. The dynamic context does not occupy an additional dimension in the nine-dimensional node feature vector. Instead, it is provided as independent graph-level state information and is fused with the embedding of the node occupied by the shuttle. This enables the agent to explicitly distinguish between the dynamic phase in which further target activations may still occur and the phase in which the dynamic activation process has terminated.

During the dynamic order process, target activation is represented by updating the state of an existing storage-location node rather than modifying the physical graph topology. Once a non-target item is activated by a newly arriving order, its task attribute changes from the inactive non-target state to an active target state with an associated due date, which is reflected in the sixth node feature. At the same time, its ninth feature is updated according to the currently adopted reward model to represent the corresponding buffered tardiness, $q$, or $R$. Hence, the static physical structure of the warehouse is primarily represented by the graph nodes and edges, whereas the dynamically evolving item status, due dates, shuttle state, current time, dynamic activation state, and lower-bound information are represented through the node-level and graph-level features.

\section{Action space}

To enable the deep reinforcement learning (DRL) agent to jointly handle the retrieval of target items and the relocation of blocking non-target items, a state-dependent upper-level action space is constructed according to the current shuttle loading state and the task attribute of the carried item. At each upper-level decision point, an action selected by the DRL agent does not directly correspond to an atomic shuttle movement. Instead, the agent selects an upper-level destination node from the currently feasible candidate nodes, while the subsequent movement is handled by the lower-level path planning and execution module.

For a state $s_t$, the upper-level action space is defined as

\[
\mathcal{A}(s_t)=
\begin{cases}
\mathcal{A}_{\mathrm{item}}(s_t), & \text{if the shuttle is unloaded},\\
\mathcal{A}_{\mathrm{I/O}}, & \text{if the shuttle carries a target item},\\
\mathcal{A}_{\mathrm{empty}}(s_t), & \text{if the shuttle carries a non-target item}.
\end{cases}
\]

where $\mathcal{A}(s_t)$ denotes the complete upper-level action space in state $s_t$, $\mathcal{A}_{\mathrm{item}}(s_t)$ denotes the set of item nodes that can be selected when the shuttle is unloaded, $\mathcal{A}_{\mathrm{I/O}}$ denotes the set of selectable I/O points, and $\mathcal{A}_{\mathrm{empty}}(s_t)$ denotes the set of empty storage locations that can be selected as relocation destinations.

When the shuttle is unloaded, the action space consists of currently directly accessible item nodes that are relevant to the scheduling process. Specifically, $\mathcal{A}_{\mathrm{item}}(s_t)$ contains two types of candidate items. The first type consists of directly accessible target items, whose selection initiates the corresponding retrieval operation. The second type consists of directly accessible non-target items that block active target items located deeper in the same storage lane. Selecting such an item initiates the relocation operation required to release the corresponding target item. Accordingly, the candidate set can be expressed conceptually as

\[
\mathcal{A}_{\mathrm{item}}(s_t)
=
\mathcal{A}_{\mathrm{target}}(s_t)
\cup
\mathcal{A}_{\mathrm{blocker}}(s_t),
\]

where $\mathcal{A}_{\mathrm{target}}(s_t)$ denotes the set of currently directly accessible target-item candidates, and $\mathcal{A}_{\mathrm{blocker}}(s_t)$ denotes the set of blocking non-target-item candidates that must be relocated to access active target items located behind them.

Here, ``directly accessible'' means that the shuttle can access the corresponding item without first moving any other occupied item. Due to the access constraints within a multi-deep storage lane, a target item blocked by other occupied items is not directly included in the current action space. Instead, the blocking non-target item located in front of it must first be processed. In contrast, if a directly accessible non-target item does not block any active target item, relocating it does not contribute to the current retrieval objective and it is therefore excluded from the current feasible action set.

In addition, a blocking non-target item constitutes a feasible candidate action only when a feasible relocation destination exists. Consequently, $\mathcal{A}_{\mathrm{item}}(s_t)$ contains only upper-level actions that are executable in the current state and contribute to the subsequent retrieval process.

When the shuttle carries a target item, the next upper-level decision is to deliver the item to an I/O point. The action space therefore becomes

\[
\mathcal{A}(s_t)=\mathcal{A}_{\mathrm{I/O}},
\]

where $\mathcal{A}_{\mathrm{I/O}}$ consists of the selectable I/O points. The DRL agent selects one of these I/O points as the delivery destination of the currently carried target item. Since different I/O points may correspond to different travel distances, I/O-point selection itself constitutes part of the retrieval scheduling decision.

When the shuttle carries a blocking non-target item, the item cannot be delivered to an I/O point and must instead be relocated to a feasible empty storage location in another storage lane. In this case, the action space becomes

\[
\mathcal{A}(s_t)=\mathcal{A}_{\mathrm{empty}}(s_t).
\]

To satisfy the compact storage lane constraint adopted in this work, the relocation destination is restricted to the deepest accessible empty location in each feasible destination lane. For a partially occupied lane, this location corresponds to the deepest empty position in front of the current occupied block that preserves contiguous storage within the lane. For a completely empty lane, it corresponds to the deepest storage location in that lane. Consequently, each feasible destination lane with remaining capacity contributes at most one relocation candidate.

Furthermore, the relocation must be performed to a storage lane different from the source lane of the blocking non-target item. Therefore, the source lane is excluded from the feasible relocation destinations for the current blocker. The relocation action space can thus be expressed as

\[
\mathcal{A}_{\mathrm{empty}}(s_t)
=
\left\{
e_l
\mid
l \neq l_{\mathrm{src}},
\;
e_l \text{ is the deepest accessible empty location in lane } l
\right\},
\]

where $e_l$ denotes the deepest accessible empty location in destination lane $l$, and $l_{\mathrm{src}}$ denotes the source storage lane of the currently carried blocking non-target item. This design satisfies the compact storage lane constraint while preserving the relocation scheduling flexibility among different destination lanes, allowing the DRL agent to select an appropriate relocation destination according to the current warehouse state.

Accordingly, a complete upper-level item-processing process follows one of two basic decision sequences depending on the item type:

\[
\text{target item}
\rightarrow
\text{I/O point},
\]

or

\[
\text{blocking non-target item}
\rightarrow
\text{empty storage location}.
\]

The former corresponds to a retrieval operation, whereas the latter corresponds to a relocation operation. Both operations consist of two consecutive upper-level decisions: the item to be processed is first selected while the shuttle is unloaded, after which the corresponding destination is selected according to the task attribute of the loaded item. In this way, retrieval and relocation can be handled within a unified upper-level node-selection framework.

Since the warehouse layout, current occupancy, active target set, blocking relationships, and shuttle loading state evolve throughout the scheduling process, the sizes of $\mathcal{A}_{\mathrm{item}}(s_t)$ and $\mathcal{A}_{\mathrm{empty}}(s_t)$ also vary dynamically with the state. Therefore, the proposed action space is inherently a variable-size action space rather than a discrete action space with a fixed number of categories. At each upper-level decision point, the DRL agent only needs to select an action from the currently feasible candidate-node set.

This state-dependent action-space design excludes physically infeasible nodes and nodes that do not contribute to the current scheduling objective, thereby focusing the policy search on meaningful retrieval and relocation decisions. Moreover, since each action corresponds to a feasible node under the current state rather than to a predefined fixed action index, the same action representation can naturally accommodate different numbers of storage lanes, different lane depths, and candidate sets that continuously change with dynamic target activation.

\section{Reward function design}
\subsection{Lower-bound derivation for simple retrieval and relocation operations}

The overall scheduling process can be decomposed into a sequence of individual retrieval and relocation operations. To derive a lower bound on the completion time of these operations, several constraints associated with the actual execution process are relaxed, while only the minimum movements and loading/unloading operations that are unavoidable for completing the corresponding operation are retained. Based on this relaxation, lower bounds on the execution times of individual retrieval and relocation operations are derived as follows.

For a target item \(i\) stored in the warehouse, a basic retrieval operation consists of the following processes: the shuttle reaches the storage location of target item \(i\) in its storage lane and loads the item, transports the loaded item to an Input/Output (I/O) point, and finally unloads the item at the I/O point. Let \(d_i\) denote the depth of the storage location occupied by target item \(i\) within its storage lane. Accordingly, \(d_i\) is also equal to the minimum travel distance from the entrance of this storage lane to the storage location of target item \(i\). To obtain an optimistic estimate of the retrieval time, the travel required for the shuttle to reach the entrance of the storage lane is relaxed, and \(d_i\) is therefore used as a lower bound on the travel distance required for the shuttle to reach the storage location of target item \(i\).

Furthermore, let \(d^{\mathrm{I/O}}\) denote the static shortest-path distance from the storage location of target item \(i\) to the nearest I/O point. This distance is used as a lower bound on the travel distance required to transport target item \(i\) from its storage location to an I/O point after loading. Let \(v^{\mathrm{shuttle}}\) denote the traveling speed of the shuttle, \(\tau^{\mathrm{load}}\) the time required for one loading operation, and \(\tau^{\mathrm{unload}}\) the time required for one unloading operation. The theoretical lower bound on the retrieval operation time can then be expressed as

\begin{equation}
\tau^{\mathrm{retrieval}}
=
\frac{d_i}{v^{\mathrm{shuttle}}}
+
\tau^{\mathrm{load}}
+
\frac{d^{\mathrm{I/O}}}{v^{\mathrm{shuttle}}}
+
\tau^{\mathrm{unload}}.
\label{eq:retrieval_lb_general}
\end{equation}

In Eq.~\eqref{eq:retrieval_lb_general}, \(d_i/v^{\mathrm{shuttle}}\) represents the minimum travel time required for the shuttle to move from the entrance of the storage lane to the storage location of target item \(i\), while \(d^{\mathrm{I/O}}/v^{\mathrm{shuttle}}\) represents the minimum travel time required to transport the loaded target item from its storage location to the nearest I/O point. The terms \(\tau^{\mathrm{load}}\) and \(\tau^{\mathrm{unload}}\) represent the time required for the loading and unloading operations, respectively.

In the environment considered in this study, the shuttle moves by one adjacent location per time unit. Therefore, its traveling speed is defined as
\[
v^{\mathrm{shuttle}} = 1.
\]

Moreover, each loading and unloading operation requires one time unit, i.e.,

\[
\tau^{\mathrm{load}}
=
\tau^{\mathrm{unload}}
=
1.
\]

By substituting these environment-specific settings into Eq.~\eqref{eq:retrieval_lb_general}, the retrieval operation time lower bound is simplified to

\begin{equation}
\boxed{
\tau^{\mathrm{retrieval}}
=
d_i
+
d^{\mathrm{I/O}}
+
2
}.
\label{eq:retrieval_lb}
\end{equation}

The constant term \(2\) in Eq.~\eqref{eq:retrieval_lb} corresponds to one loading operation and one unloading operation. It should be noted that this lower bound is not intended to predict the actual completion time of the retrieval operation. Instead, it provides an optimistic estimate by neglecting additional travel that may be required before the shuttle reaches the entrance of the corresponding storage lane, waiting time caused by other tasks, and other factors that may increase the actual retrieval time. Consequently, \(\tau^{\mathrm{retrieval}}\) represents a lower bound on the time required to complete the retrieval operation rather than an estimate of its actual execution time.

The formulation in Eq.~\eqref{eq:retrieval_lb} corresponds to the general case in which target item \(i\) is still stored at its original storage location and has not yet been loaded. The lower-bound calculation is further adapted to intermediate system states. If the shuttle is already located at the storage location of target item \(i\), the approach movement represented by \(d_i\) has already been completed and is therefore excluded from the remaining operation-time lower bound. Furthermore, if target item \(i\) has already been loaded onto the shuttle, both the approach movement and loading operation have already been completed. In this case, only the minimum travel to the nearest I/O point and the final unloading operation remain. These state-dependent adjustments provide a tighter estimate of the remaining retrieval time without changing the underlying relaxation principle.

For a relocation operation, consider a blocking non-target item \(j\) that prevents the retrieval of a target item located deeper in the same storage lane. Let \(d_j\) denote the depth of the storage location occupied by blocking non-target item \(j\) within its original storage lane. To complete the relocation operation, the shuttle must at least perform the following processes: it first travels from the entrance of the original storage lane to the storage location of blocking non-target item \(j\) and loads the item; it then transports the loaded item out of the original storage lane; subsequently, it enters a feasible destination storage lane; and finally, it unloads the blocking item at a feasible storage location in the destination lane.

Similar to the retrieval operation, the actual relocation process is relaxed to obtain a lower bound that does not overestimate the required operation time. In particular, the actual relocation destination of blocking non-target item \(j\) is not explicitly determined when calculating the relocation operation time lower bound. Instead, it is assumed that a feasible destination storage lane with the minimum possible entry cost is available. Moreover, the additional cross-lane travel distance between the entrance of the original storage lane and the entrance of the actual destination storage lane is neglected. Let \(d^{\min,\mathrm{entry}}\) denote the minimum unavoidable travel distance required for the shuttle carrying the blocking item to enter a feasible destination storage lane. The theoretical lower bound on the relocation operation time is then expressed as

\begin{equation}
\tau^{\mathrm{relocation}}
=
\frac{d_j}{v^{\mathrm{shuttle}}}
+
\tau^{\mathrm{load}}
+
\frac{d_j}{v^{\mathrm{shuttle}}}
+
\frac{d^{\min,\mathrm{entry}}}{v^{\mathrm{shuttle}}}
+
\tau^{\mathrm{unload}}.
\label{eq:relocation_lb_general}
\end{equation}

In Eq.~\eqref{eq:relocation_lb_general}, the first term \(d_j/v^{\mathrm{shuttle}}\) represents the minimum travel time required for the shuttle to approach blocking non-target item \(j\) from the entrance of its original storage lane. The second term \(d_j/v^{\mathrm{shuttle}}\) represents the minimum travel time required to transport the loaded blocking item completely out of its original storage lane. The term \(d^{\min,\mathrm{entry}}/v^{\mathrm{shuttle}}\) represents the minimum unavoidable travel time required to enter a feasible destination storage lane. Finally, \(\tau^{\mathrm{load}}\) and \(\tau^{\mathrm{unload}}\) correspond to the loading and unloading operations, respectively. Since the cross-lane travel distance between the original storage lane and a specific destination storage lane is completely relaxed, this distance is not included in the lower-bound calculation.

Under the environment settings considered in this study, \(v^{\mathrm{shuttle}}=1\), and both loading and unloading require one time unit. In addition, under the relaxed relocation model, entering a feasible destination storage lane requires at least one movement to an adjacent location. Therefore,
\[
d^{\min,\mathrm{entry}} = 1.
\]

By substituting

\[
v^{\mathrm{shuttle}}=1,
\qquad
\tau^{\mathrm{load}}=1,
\qquad
\tau^{\mathrm{unload}}=1,
\qquad
d^{\min,\mathrm{entry}}=1
\]

into Eq.~\eqref{eq:relocation_lb_general}, the relocation operation time lower bound is simplified to

\begin{equation}
\boxed{
\tau^{\mathrm{relocation}}
=
2d_j
+
3
}.
\label{eq:relocation_lb}
\end{equation}

In Eq.~\eqref{eq:relocation_lb}, the term \(2d_j\) accounts for the minimum travel required to first approach blocking non-target item \(j\) and subsequently transport it out of its original storage lane. The constant term \(3\) consists of one loading operation, one minimum movement required to enter a feasible destination storage lane, and one unloading operation.

Similar to the retrieval operation time lower bound, Eq.~\eqref{eq:relocation_lb} does not determine the actual relocation destination of blocking non-target item \(j\), nor does it include the additional travel distance that may be required between the original storage lane and the actual destination lane. Hence, \(\tau^{\mathrm{relocation}}\) represents an optimistic lower bound on the time required to complete the relocation operation.

The lower-bound calculation is also adapted when blocking non-target item \(j\) has already been loaded onto the shuttle. In such an intermediate state, the approach movement toward item \(j\) and its loading operation have already been completed and are therefore excluded from the remaining relocation time lower bound. Only the unavoidable remaining movements required to leave the original storage lane, enter a feasible relocation lane, and unload the item are retained.

Through the above relaxations, a minimum operation time can be derived for both the retrieval of each target item and the relocation of each blocking non-target item. These individual operation-time lower bounds subsequently serve as the basic components for constructing the state-level lower bound. In the following derivation, the unavoidable relocation work caused by blocking non-target items and the precedence relationships among target items imposed by the multi-deep storage-lane topology are further incorporated to obtain a lower bound on the total tardiness of the current system state.

\subsection{Lower-bound derivation within a single storage lane}

In the previous subsection, the retrieval operation time lower bound
\(\tau^{\mathrm{retrieval}}\) for an individual target item and the relocation
operation time lower bound \(\tau^{\mathrm{relocation}}\) for an individual
blocking non-target item were derived by relaxing part of the movement and
scheduling constraints involved in the actual execution process. However, in a
multi-deep storage lane, the retrieval operation of a target item generally
cannot be completed independently.

Although different scheduling policies may generate different global operation
sequences over the entire warehouse layout, the necessary access precedence
among the items currently stored within a given multi-deep storage lane is
determined by its Last-In, First-Out-type (LIFO-type) access constraint and is
therefore independent of the upper-level scheduling policy. Specifically,
because a storage lane can only be accessed from its lane entrance, an item
located deeper in the lane cannot be retrieved or relocated before the items
blocking access to it have been removed. Consequently, although different
scheduling policies may change the operation sequence across different storage
lanes as well as the actual execution time of individual operations, the
necessary in-lane access precedence induced by the current depth ordering of
the items remains unchanged.

Based on this property, the local precedence structure within each storage lane
is further exploited to aggregate the individual retrieval and relocation
operation time lower bounds derived in the previous subsection and construct
lane-level lower bounds. Specifically, a target-specific required-work lower
bound \(q\) is first constructed for each target item to represent the minimum
amount of work required to retrieve the target item together with the necessary
relocations of its preceding blockers. A precedence-constrained cumulative work
lower bound \(R\) is subsequently constructed to represent the minimum
cumulative work required to complete a given number of target items while
respecting the necessary in-lane access precedence.

Consider an arbitrary storage lane \(\ell\). All occupied storage locations in
this lane are scanned from the lane entrance toward the deeper end of the lane.
The target items encountered during this scan are indexed according to their
scanning order as

\[
i^{\mathrm{target}}_{\ell,1},
\;
i^{\mathrm{target}}_{\ell,2},
\;
\ldots,
\;
i^{\mathrm{target}}_{\ell,n_\ell},
\]

where \(\ell\) denotes the index of the storage lane, \(n_\ell\) denotes the
number of target items currently contained in storage lane \(\ell\), and
\(i^{\mathrm{target}}_{\ell,r}\) denotes the \(r\)-th target item encountered
when scanning from the lane entrance toward the deeper end of the lane.
Accordingly, a smaller value of \(r\) indicates a target item located closer to
the lane entrance, whereas a larger value of \(r\) indicates a target item
located deeper in the storage lane.

During this scan, if one or more non-target items are encountered before the
next target item, these non-target items are temporarily recorded as blockers.
When the next target item \(i^{\mathrm{target}}_{\ell,r}\) is encountered, all
blockers that have not yet been assigned are associated with this target item.
Let

\[
\mathcal{I}^{\mathrm{blocker}}_{\ell,r}
\]

denote the set of blocking non-target items associated with target item
\(i^{\mathrm{target}}_{\ell,r}\). The \(b\)-th blocking non-target item in
storage lane \(\ell\) is denoted by

\[
i^{\mathrm{blocker}}_{\ell,b}.
\]

Thus, \(\mathcal{I}^{\mathrm{blocker}}_{\ell,r}\) contains the blocking
non-target items that are located in front of
\(i^{\mathrm{target}}_{\ell,r}\), must be relocated before this target item can
be retrieved, and have not already been associated with another target item
closer to the lane entrance.

For example, consider a local item arrangement in storage lane \(\ell\), from
the lane entrance toward the deeper end, given by

\[
i^{\mathrm{blocker}}_{\ell,1},
\;
i^{\mathrm{blocker}}_{\ell,2},
\;
i^{\mathrm{target}}_{\ell,1},
\;
i^{\mathrm{blocker}}_{\ell,3},
\;
i^{\mathrm{target}}_{\ell,2}.
\]

The corresponding blocker sets are

\[
\mathcal{I}^{\mathrm{blocker}}_{\ell,1}
=
\left\{
i^{\mathrm{blocker}}_{\ell,1},
i^{\mathrm{blocker}}_{\ell,2}
\right\}
\]

and

\[
\mathcal{I}^{\mathrm{blocker}}_{\ell,2}
=
\left\{
i^{\mathrm{blocker}}_{\ell,3}
\right\}.
\]

This blocker-to-target association ensures that each blocking non-target item
is counted only once when constructing the required-work lower bounds, thereby
preventing the same relocation operation from being repeatedly included in the
lower-bound calculations of multiple target items.

It should be noted that if a non-target item is located deeper than the last
target item currently contained in the storage lane, and no target item exists
behind it, this non-target item is not included in any
\(\mathcal{I}^{\mathrm{blocker}}_{\ell,r}\). This is because relocating such an
item is not necessarily required to retrieve any of the target items currently
contained in the lane. Including such unnecessary relocation work could
overestimate the minimum required work and would therefore violate the
lower-bound property.

\subsubsection{Target-specific required-work lower bound \(q\).}

For a target item \(i^{\mathrm{target}}_{\ell,r}\) in storage lane \(\ell\),
its retrieval under the current storage configuration requires at least two
types of unavoidable work. First, the retrieval operation of
\(i^{\mathrm{target}}_{\ell,r}\) itself must be completed. Second, if blocking
non-target items exist in front of the target item, all items in
\(\mathcal{I}^{\mathrm{blocker}}_{\ell,r}\) must first be relocated.

Accordingly, the target-specific required-work lower bound associated with
target item \(i^{\mathrm{target}}_{\ell,r}\) is defined as

\begin{equation}
q_{\ell,r}
=
\tau^{\mathrm{retrieval}}
\left(
i^{\mathrm{target}}_{\ell,r}
\right)
+
\sum_{
i^{\mathrm{blocker}}_{\ell,b}
\in
\mathcal{I}^{\mathrm{blocker}}_{\ell,r}
}
\tau^{\mathrm{relocation}}
\left(
i^{\mathrm{blocker}}_{\ell,b}
\right).
\label{eq:target_specific_required_work_lb}
\end{equation}

In Eq.~\eqref{eq:target_specific_required_work_lb}, \(q_{\ell,r}\) denotes the
minimum amount of required work associated with the \(r\)-th target item in
storage lane \(\ell\);
\(\tau^{\mathrm{retrieval}}
(i^{\mathrm{target}}_{\ell,r})\) denotes the retrieval operation time lower
bound of target item \(i^{\mathrm{target}}_{\ell,r}\);
\(\mathcal{I}^{\mathrm{blocker}}_{\ell,r}\) denotes the set of blocking
non-target items associated with this target item; and
\(\tau^{\mathrm{relocation}}
(i^{\mathrm{blocker}}_{\ell,b})\) denotes the relocation operation time lower
bound of blocking non-target item \(i^{\mathrm{blocker}}_{\ell,b}\).

If no blocking non-target item needs to be relocated before
\(i^{\mathrm{target}}_{\ell,r}\) can be retrieved, i.e.,

\[
\mathcal{I}^{\mathrm{blocker}}_{\ell,r}
=
\varnothing,
\]

the target-specific required-work lower bound reduces to

\[
q_{\ell,r}
=
\tau^{\mathrm{retrieval}}
\left(
i^{\mathrm{target}}_{\ell,r}
\right).
\]

Hence, \(q_{\ell,r}\) represents the minimum local amount of work that is
unavoidably required to complete target item
\(i^{\mathrm{target}}_{\ell,r}\), while temporarily disregarding the
cumulative precedence effect of the other target items in the same storage
lane. Compared with the retrieval operation time lower bound alone,
\(q_{\ell,r}\) explicitly incorporates the necessary relocation work induced
by the blocking relationship in front of the corresponding target item.

Furthermore, both \(\tau^{\mathrm{retrieval}}\) and
\(\tau^{\mathrm{relocation}}\) are evaluated according to the current system
state. Therefore, if the shuttle is already located at the location of a
particular item, or if a target item or blocking non-target item has already
been loaded onto the shuttle, the movement or loading operations that have
already been completed are not counted again in the corresponding operation
time lower bound. Consequently, \(q_{\ell,r}\) represents the minimum
remaining amount of work required from the current state to complete the tasks
associated with target item \(i^{\mathrm{target}}_{\ell,r}\), rather than the
complete operation time recalculated from its initial storage state.

\subsubsection{Precedence-constrained cumulative work lower bound \(R\).}

Although \(q_{\ell,r}\) incorporates the necessary work associated with an
individual target item and its direct blockers, it does not by itself fully
represent the precedence relationship among different target items induced by
the LIFO-type access constraint of a multi-deep storage lane.

Consider two target items in storage lane \(\ell\),

\[
i^{\mathrm{target}}_{\ell,r-1}
\quad \text{and} \quad
i^{\mathrm{target}}_{\ell,r},
\]

ordered from the lane entrance toward the deeper end. Since
\(i^{\mathrm{target}}_{\ell,r}\) is located deeper in the lane, its retrieval
cannot be completed while bypassing the necessary preceding work associated
with the items located closer to the lane entrance. More generally, to
complete the \(r\)-th target item in storage lane \(\ell\), all unavoidable
target-related work associated with the preceding part of the lane forms
necessary preceding work.

Therefore, after the target-specific required-work lower bounds
\(q_{\ell,r}\) have been obtained, they are accumulated according to the
target order from the lane entrance toward the deeper end. The
precedence-constrained cumulative work lower bound is defined as

\begin{equation}
R_{\ell,r}
=
\sum_{h=1}^{r}
q_{\ell,h}.
\label{eq:precedence_cumulative_work_lb}
\end{equation}

In Eq.~\eqref{eq:precedence_cumulative_work_lb}, \(R_{\ell,r}\) denotes the
minimum cumulative amount of work required, from the current state, to complete
the first \(r\) target items of storage lane \(\ell\) while respecting the
necessary in-lane access precedence, and \(q_{\ell,h}\) denotes the
target-specific required-work lower bound associated with the \(h\)-th target
item in this storage lane.

Accordingly,
\[
R_{\ell,1}
=
q_{\ell,1},
\]

\[
R_{\ell,2}
=
q_{\ell,1}
+
q_{\ell,2},
\]

and, in general,

\begin{equation}
R_{\ell,r}
=
R_{\ell,r-1}
+
q_{\ell,r}.
\end{equation}

In addition,

\[
R_{\ell,0}=0
\]

is defined to represent the case in which no target item is completed from
storage lane \(\ell\), such that this lane contributes no additional required
work.

Consequently, for storage lane \(\ell\) containing \(n_\ell\) current target
items, a sequence of precedence-constrained cumulative work lower bounds is
obtained as

\[
R_{\ell,1},
\;
R_{\ell,2},
\;
\ldots,
\;
R_{\ell,n_\ell}.
\]

Each \(R_{\ell,r}\) represents the minimum cumulative work required to complete
the first \(r\) target items in storage lane \(\ell\) while preserving the
necessary in-lane precedence relation.

From the perspective of lower-bound construction, \(q\) and \(R\) describe
two different levels of information. The target-specific required-work lower
bound \(q_{\ell,r}\) captures the minimum local work associated with target item
\(i^{\mathrm{target}}_{\ell,r}\) and its direct blocking non-target items,
whereas the precedence-constrained cumulative work lower bound \(R_{\ell,r}\)
further incorporates the target-item precedence induced by the LIFO-type access
constraint of the storage lane. Therefore, \(R_{\ell,r}\) includes not only
the required work associated with target item \(i^{\mathrm{target}}_{\ell,r}\)
itself, but also all preceding target-related work that cannot be bypassed
before this deeper target item can be completed.

For example, consider three target items in storage lane \(\ell\), ordered from
the lane entrance toward the deeper end as
\[
i^{\mathrm{target}}_{\ell,1},
\quad
i^{\mathrm{target}}_{\ell,2},
\quad
i^{\mathrm{target}}_{\ell,3},
\]

with the corresponding target-specific required-work lower bounds
\[
q_{\ell,1},
\quad
q_{\ell,2},
\quad
q_{\ell,3}.
\]

The corresponding precedence-constrained cumulative work lower bounds are

\[
R_{\ell,1}
=
q_{\ell,1},
\]

\[
R_{\ell,2}
=
q_{\ell,1}
+
q_{\ell,2},
\]

and

\[
R_{\ell,3}
=
q_{\ell,1}
+
q_{\ell,2}
+
q_{\ell,3}.
\]

Thus, if the third and deeper target item is to be completed, its minimum
required work cannot be characterized solely by \(q_{\ell,3}\). The
unavoidable target-related work associated with the items located in front of
it must also be considered. Therefore, \(R_{\ell,3}\) provides a tighter
cumulative work lower bound while preserving the precedence structure of the
current storage lane.

Through the above construction, the operation-level lower bounds

\[
\tau^{\mathrm{retrieval}}
\quad \text{and} \quad
\tau^{\mathrm{relocation}}
\]

are first aggregated into the target-level target-specific required-work lower
bound \(q_{\ell,r}\), which is subsequently accumulated according to the
deterministic in-lane precedence relation to obtain the
precedence-constrained cumulative work lower bound \(R_{\ell,r}\). The
single-lane lower-bound construction can therefore be summarized as

\[
\left\{
\tau^{\mathrm{retrieval}},
\tau^{\mathrm{relocation}}
\right\}
\longrightarrow
q_{\ell,r}
\longrightarrow
R_{\ell,r}.
\]

At this stage, each storage lane is represented by a set of
precedence-constrained cumulative work lower bounds. In the subsequent
derivation, the \(R_{\ell,r}\) values across different storage lanes are
combined while allowing the most favorable combination of completed target
items among the lanes. This provides global completion-time lower bounds for
different numbers of completed target items and ultimately enables the
construction of a lower bound on the total tardiness of the current warehouse
state.

\subsection{Layout-level lower-bound derivation}

In the previous subsection, the necessary access precedence induced by the
Last-In, First-Out-type (LIFO-type) access constraint within an individual
storage lane was exploited to first aggregate the lower bounds on individual
retrieval and relocation operation times into the target-specific
required-work lower bound \(q_{\ell,r}\), and subsequently construct the
precedence-constrained cumulative work lower bound \(R_{\ell,r}\). Specifically,

\[
q_{\ell,r}
=
\tau^{\mathrm{retrieval}}
\left(
i^{\mathrm{target}}_{\ell,r}
\right)
+
\sum_{
i^{\mathrm{blocker}}_{\ell,b}
\in
\mathcal{I}^{\mathrm{blocker}}_{\ell,r}
}
\tau^{\mathrm{relocation}}
\left(
i^{\mathrm{blocker}}_{\ell,b}
\right)
\]

represents the minimum required work directly associated with target item
\(i^{\mathrm{target}}_{\ell,r}\), whereas

\[
R_{\ell,r}
=
\sum_{h=1}^{r}
q_{\ell,h}
\]

further represents the minimum cumulative work that is unavoidable for
completing the first \(r\) target items in storage lane \(\ell\) while
respecting the necessary in-lane access precedence.

However, \(q_{\ell,r}\) and \(R_{\ell,r}\) only characterize the local work
structure within an individual storage lane, whereas the complete warehouse
layout generally contains multiple storage lanes with target items to be
processed. Unlike the necessary access precedence within the same storage lane,
which is determined by the LIFO-type access constraint, the processing order
among different storage lanes is not uniquely determined by the current
physical layout and may vary depending on the scheduling policy. Therefore, to
construct a layout-level lower bound that remains valid for any feasible
scheduling policy, the specific retrieval sequence of the target items is not
predefined at the complete-layout level, nor is a particular completion rank
associated with a specific target item. Instead, the derivation focuses on the
minimum cumulative amount of work that is unavoidable for completing a given
number of target items over all possible scheduling sequences.

To this end, let \(n^{\mathrm{target}}\) denote the number of currently
unfinished target items in the warehouse layout. For each completion rank

\[
k
=
1,2,\ldots,n^{\mathrm{target}},
\]

\(C_k\) is defined as the cumulative work lower bound that is unavoidable for
completing any \(k\) of the current target items from the current system state
under the corresponding constraint relaxation. Here, \(k\) represents only the
number of completed target items and is not associated with any predefined
specific target item. Therefore, \(C_k\) does not represent the operation-time
lower bound of the \(k\)-th specific retrieval task. Instead, it represents the
minimum cumulative amount of work that any feasible schedule must have
performed before the \(k\)-th target-item completion can occur. In other words,
\(C_k\) is a cumulative work lower bound associated with completion rank \(k\),
rather than the processing time of a particular target item.

Based on this common definition, three layout-level lower-bound calculation
modes are proposed, namely \texttt{simple\_shuttle\_LB},
\texttt{mixed\_shuttle\_LB}, and \texttt{topology\_shuttle\_LB}. All three
modes employ the same lower bounds on individual retrieval and relocation
operation times and share the same construction of the target-specific
required-work lower bound \(q_{\ell,r}\) and the precedence-constrained
cumulative work lower bound \(R_{\ell,r}\). In addition, all three modes
preserve the serial workload constraint imposed by the use of a single
shuttle. That is, the work required by multiple retrieval-related tasks cannot
be processed in parallel by multiple shuttles, but must instead be performed
sequentially by the same shuttle. The fundamental difference among the three
modes lies in the extent to which the precedence information induced by the
LIFO-type access constraints of the individual storage lanes is preserved or
relaxed when the single-lane lower bounds are aggregated into the layout-level
cumulative work lower bounds \(C_k\).

Specifically, \texttt{simple\_shuttle\_LB} applies the strongest relaxation to
the precedence relations among target items within the same storage lane. It
retains only the target-specific required work and the serial workload
accumulation imposed by the single shuttle. Building on these constraints,
\texttt{mixed\_shuttle\_LB} additionally exploits the lane-level precedence
information contained in \(R_{\ell,r}\) to tighten the corresponding cumulative
work lower bounds, while still relaxing the complete combinatorial consistency
between different precedence-constrained cumulative work lower bounds and the
storage lanes to which they belong. Finally,
\texttt{topology\_shuttle\_LB} preserves the complete
precedence-constrained cumulative work structure induced by the LIFO-type
access constraint within each storage lane and only allows the most favorable
combination to be selected across different storage lanes for completing a
given number of target items. Consequently, the three modes form a progressive
layout-level lower-bound construction framework ranging from a stronger
constraint relaxation to a more complete preservation of the storage-lane
topology, yielding

\[
C_k^{\mathrm{simple}},
\qquad
C_k^{\mathrm{mixed}},
\qquad
C_k^{\mathrm{topology}}.
\]

\subsubsection{Simple shuttle lower-bound mode.}

The \texttt{simple\_shuttle\_LB} mode retains two
fundamental constraints. First, for each target item
\(i^{\mathrm{target}}_{\ell,r}\), neither its own retrieval operation nor the
necessary relocation operations of its directly associated blocking non-target
items can be omitted. Therefore, the corresponding target-specific
required-work lower bound \(q_{\ell,r}\) is fully retained. Second, because all
retrieval-related operations in the considered system are executed by a single
shuttle, the required work associated with different target items cannot be
processed in parallel by multiple shuttles and must instead be performed
serially by the same shuttle. Consequently, the minimum required work
associated with multiple completed target items must be accumulated.

At the same time, \texttt{simple\_shuttle\_LB} completely relaxes the
precedence relations among different target items within the same storage lane
that are induced by the Last-In, First-Out-type (LIFO-type) access constraint.
Under this relaxation, all target-specific required-work lower bounds
\(q_{\ell,r}\) in the entire layout are treated as independent required-work
units. Neither the depth position of a target item relative to other target
items in the same storage lane nor the target-related work located closer to
the lane entrance is required to be considered before that target item is
selected. Accordingly, the precedence-constrained cumulative structure
contained in \(R_{\ell,r}\) is not used to restrict the target-item completion
order in this mode.

Based on this relaxation, all target-specific required-work lower bounds
associated with the unfinished target items in the current warehouse layout
are collected into the set

\[
\mathcal{Q}
=
\left\{
q_{\ell,r}
\mid
i^{\mathrm{target}}_{\ell,r}
\text{ is an unfinished target item}
\right\}.
\]

Let \(n^{\mathrm{target}}\) denote the number of unfinished target items in the
current layout. Accordingly, \(\mathcal{Q}\) contains
\(n^{\mathrm{target}}\) target-specific required-work lower bounds. These
values are sorted in ascending order without considering their original
storage-lane membership or their in-lane precedence positions, yielding

\[
q^{\mathrm{sorted}}_{1}
\leq
q^{\mathrm{sorted}}_{2}
\leq
\cdots
\leq
q^{\mathrm{sorted}}_{n^{\mathrm{target}}}.
\]

Here, \(q^{\mathrm{sorted}}_{h}\) denotes the \(h\)-th smallest
target-specific required-work lower bound in the entire current layout. This
sorting is performed solely according to the magnitude of the required-work
lower bounds and does not preserve the original storage-lane membership or
target-item precedence information. Therefore, under the simple-mode
relaxation, a target item located deeper in a storage lane may be considered as
an earlier completion candidate if its \(q_{\ell,r}\) is sufficiently small,
without first enforcing the precedence relations induced by other target items
located closer to the lane entrance.

For each completion rank

\[
k
=
1,2,\ldots,n^{\mathrm{target}},
\]

the most favorable assumption is adopted to obtain the minimum cumulative work
required for completing any \(k\) target items. Specifically, the shuttle is
assumed to be able to process the \(k\) smallest target-specific required-work
units in the entire layout first. The cumulative work lower bound under
\texttt{simple\_shuttle\_LB} is therefore defined as

\begin{equation}
C_k^{\mathrm{simple}}
=
\sum_{h=1}^{k}
q^{\mathrm{sorted}}_{h}.   
\end{equation}

Here, \(C_k^{\mathrm{simple}}\) denotes the minimum cumulative amount of work
that the single shuttle must perform before any \(k\) of the current target
items can be completed under the simple relaxation;
\(q^{\mathrm{sorted}}_{h}\) denotes the \(h\)-th smallest target-specific
required-work lower bound; and \(k\) denotes the corresponding completion
rank.

Accordingly, for the first completion rank,

\[
C_1^{\mathrm{simple}}
=
q^{\mathrm{sorted}}_1,
\]

and for the second completion rank,

\[
C_2^{\mathrm{simple}}
=
q^{\mathrm{sorted}}_1
+
q^{\mathrm{sorted}}_2.
\]

More generally,

\[
C_k^{\mathrm{simple}}
=
C_{k-1}^{\mathrm{simple}}
+
q^{\mathrm{sorted}}_k.
\]

This accumulation reflects the single-shuttle serial workload constraint.
Even under the most favorable scheduling conditions, if \(k\) target items are
to be completed, the required work associated with these target items cannot
be executed simultaneously by multiple parallel resources. Instead, it must
be performed sequentially by the same shuttle. Hence, the minimum cumulative
work required before the \(k\)-th target-item completion cannot be smaller than
the sum of the required-work lower bounds associated with the selected \(k\)
target items.

Furthermore, because the simple mode does not prescribe a specific retrieval
sequence, the most optimistic combination of \(k\) target items must be
considered in order to preserve the lower-bound property for any feasible
scheduling policy. Globally sorting all \(q_{\ell,r}\) values and selecting the
\(k\) smallest ones corresponds exactly to the minimum cumulative required
work attainable after the target-to-target in-lane precedence relations have
been completely relaxed. Therefore, for any feasible schedule that completes
\(k\) target items, the cumulative work actually performed by the single
shuttle cannot be smaller than \(C_k^{\mathrm{simple}}\).

From the perspective of constraint relaxation,
\texttt{simple\_shuttle\_LB} can therefore be characterized as retaining the
unavoidable required work associated with each target item and its direct
blockers, as well as the serial workload accumulation imposed by the single
shuttle, while completely relaxing the target-to-target precedence constraints
induced by the storage-lane topology. Due to this relatively strong
relaxation, \(C_k^{\mathrm{simple}}\) generally provides the most optimistic
completion-rank cumulative-work estimate among the three proposed modes and
serves as the basis on which \texttt{mixed\_shuttle\_LB} and
\texttt{topology\_shuttle\_LB} subsequently introduce additional
storage-lane precedence information.

\subsubsection{Mixed shuttle lower-bound mode.}

The \texttt{mixed\_shuttle\_LB} mode builds upon the single-shuttle
serial-workload lower bound constructed by \texttt{simple\_shuttle\_LB} and
further introduces the precedence-constrained cumulative work lower bound
\(R_{\ell,r}\) within an individual storage lane to provide an additional
lane-internal precedence constraint for each completion rank. Its purpose is
not to replace the single-shuttle constraint retained by the simple mode with
\(R_{\ell,r}\), but rather to combine two independent necessary constraints,
namely the single-shuttle serial workload constraint and the storage-lane
precedence constraint, thereby obtaining a tighter cumulative work lower bound
than that provided by the simple mode.

First, \texttt{simple\_shuttle\_LB} derives

\[
C_k^{\mathrm{simple}}
=
\sum_{h=1}^{k}
q_h^{\mathrm{sorted}},
\]

where \(C_k^{\mathrm{simple}}\) represents the minimum cumulative amount of
work that the single shuttle must perform before any \(k\) of the current
target items can be completed when the lane-level precedence relations among
different target items are completely relaxed. Therefore,
\(C_k^{\mathrm{simple}}\) mainly preserves two fundamental constraints: the
unavoidable required work associated with each target item and its directly
associated blockers, and the shared-resource constraint that all such required
work must be performed serially by the same shuttle.

However, the simple mode treats all \(q_{\ell,r}\) values in the entire layout
as independent required-work units and allows them to be globally sorted
without preserving the target-to-target precedence relations in their original
storage lanes. Consequently, for a target item located deeper in a storage
lane, the simple mode may regard it as an early completion candidate solely
because its own \(q_{\ell,r}\) is relatively small, without further requiring
the target-related work located closer to the lane entrance to have already
been completed. To complement this relaxation, the mixed mode additionally
introduces the precedence-constrained cumulative work lower bound
\(R_{\ell,r}\).

For the \(r\)-th target item \(i^{\mathrm{target}}_{\ell,r}\) in storage lane
\(\ell\),

\[
R_{\ell,r}
=
\sum_{h=1}^{r}
q_{\ell,h}.
\]

Here, \(R_{\ell,r}\) represents the minimum cumulative amount of work that is
unavoidable for completing the first \(r\) target items in storage lane
\(\ell\) while respecting the necessary access precedence induced by the
Last-In, First-Out-type (LIFO-type) access constraint. Therefore, unlike
\(q_{\ell,r}\), which only represents the local required work associated with
target item \(i^{\mathrm{target}}_{\ell,r}\) and its directly associated
blockers, \(R_{\ell,r}\) additionally incorporates all preceding
target-related work that cannot be bypassed before the deeper target item can
be completed.

It is important to note that \(R_{\ell,r}\) is a lane-local lower bound. Its
construction only characterizes the necessary cumulative work within storage
lane \(\ell\) and does not encode the cross-lane shared-resource coupling
arising from multiple storage lanes competing for the same shuttle.
Furthermore, \(R_{\ell,r}\) inherits the relaxations already introduced in the
lower bounds on individual retrieval and relocation operation times. For
example, actual cross-lane repositioning, the specific relocation destination,
cross-aisle congestion, route conflicts, and other factors that may increase
the actual execution time are not fully considered. Therefore,
\(R_{\ell,r}\) itself is not a whole-layout single-shuttle cumulative-work
bound, but rather a precedence-constrained local-work bound for an individual
storage lane.

At the layout level, the mixed mode collects the \(R_{\ell,r}\) values
associated with all currently unfinished target items into the set

\[
\mathcal{R}
=
\left\{
R_{\ell,r}
\mid
i^{\mathrm{target}}_{\ell,r}
\text{ is an unfinished target item}
\right\}.
\]

Let \(n^{\mathrm{target}}\) denote the number of unfinished target items in the
current layout. All \(R_{\ell,r}\) values in \(\mathcal{R}\) are then sorted in
ascending order as

\[
R_1^{\mathrm{sorted}}
\leq
R_2^{\mathrm{sorted}}
\leq
\cdots
\leq
R_{n^{\mathrm{target}}}^{\mathrm{sorted}}.
\]

Here, \(R_k^{\mathrm{sorted}}\) denotes the \(k\)-th smallest
precedence-constrained cumulative work lower bound in the entire layout. This
global sorting no longer preserves the complete shared-resource coupling among
the storage lanes to which the individual \(R_{\ell,r}\) values belong. From a
relaxation perspective, it can therefore be interpreted as a highly optimistic
lane-parallel representation: the local precedence structure within each
storage lane is preserved independently, whereas the resource competition
among different storage lanes caused by sharing the same shuttle is relaxed in
this \(R\)-based bound. This interpretation does not imply that the actual
system or the lower-bound implementation assumes the physical existence of
multiple shuttles; rather, it describes the mathematical relaxation underlying
the \(R\)-based bound.

Nevertheless, \(R_k^{\mathrm{sorted}}\) remains a valid lower bound for the
\(k\)-th completion rank. Consider any feasible schedule and let

\[
\mathcal{I}_k^{\mathrm{completed}}
\]

denote the set of target items that have been completed when the \(k\)-th
target-item completion occurs, with

\[
\left|
\mathcal{I}_k^{\mathrm{completed}}
\right|
=
k.
\]

For any target item
\(i^{\mathrm{target}}_{\ell,r}
\in
\mathcal{I}_k^{\mathrm{completed}}\),
the unavoidable preceding work represented by \(R_{\ell,r}\) within its
storage lane must already have been completed before that target item can be
completed. Therefore, the actual cumulative work performed before the
\(k\)-th completion must satisfy

\[
C_k^{\mathrm{actual}}
\geq
R_{\ell,r},
\qquad
\forall\,
i^{\mathrm{target}}_{\ell,r}
\in
\mathcal{I}_k^{\mathrm{completed}}.
\]

It follows that

\[
C_k^{\mathrm{actual}}
\geq
\max_{
i^{\mathrm{target}}_{\ell,r}
\in
\mathcal{I}_k^{\mathrm{completed}}
}
R_{\ell,r}.
\]

Moreover, for any subset of \(k\) values selected from all
\(R_{\ell,r}\) values in the current layout, the maximum value of this subset
cannot be smaller than the globally \(k\)-th smallest value
\(R_k^{\mathrm{sorted}}\). Hence,

\[
\max_{
i^{\mathrm{target}}_{\ell,r}
\in
\mathcal{I}_k^{\mathrm{completed}}
}
R_{\ell,r}
\geq
R_k^{\mathrm{sorted}},
\]

which yields

\[
C_k^{\mathrm{actual}}
\geq
R_k^{\mathrm{sorted}}.
\]

Therefore, \(R_k^{\mathrm{sorted}}\) provides a completion-rank lower bound
derived from the storage-lane precedence structure without prescribing a
specific retrieval sequence.

However, this \(R\)-based lower bound does not replace the single-shuttle
serial workload constraint retained by \texttt{simple\_shuttle\_LB}. Since the
actual system contains only one shuttle, any feasible schedule must also
satisfy

\[
C_k^{\mathrm{actual}}
\geq
C_k^{\mathrm{simple}}.
\]

Consequently, for every completion rank \(k\), the actual cumulative work must
simultaneously satisfy the two independent lower-bound constraints

\[
C_k^{\mathrm{actual}}
\geq
C_k^{\mathrm{simple}}
\]

and

\[
C_k^{\mathrm{actual}}
\geq
R_k^{\mathrm{sorted}}.
\]

Based on these two conditions, \texttt{mixed\_shuttle\_LB} combines the
single-shuttle serial-workload lower bound provided by
\texttt{simple\_shuttle\_LB} with the \(R\)-based lane-precedence lower bound
and defines

\begin{equation}
C_k^{\mathrm{mixed}}
=
\max
\left\{
C_k^{\mathrm{simple}},
R_k^{\mathrm{sorted}}
\right\}.    
\end{equation}

Here, \(C_k^{\mathrm{mixed}}\) denotes the cumulative work lower bound that is
unavoidable before the \(k\)-th target-item completion under the mixed mode;
\(C_k^{\mathrm{simple}}\) denotes the lower bound derived from the
target-specific required work and the single-shuttle serial workload
constraint; and \(R_k^{\mathrm{sorted}}\) denotes the additional lower-bound
constraint derived from the storage-lane precedence information.

Because both \(C_k^{\mathrm{simple}}\) and \(R_k^{\mathrm{sorted}}\) are valid
lower bounds on \(C_k^{\mathrm{actual}}\), taking their maximum also preserves
the lower-bound property. By definition,

\[
C_k^{\mathrm{mixed}}
\geq
C_k^{\mathrm{simple}}.
\]

If

\[
R_k^{\mathrm{sorted}}
\leq
C_k^{\mathrm{simple}},
\]

the lane-precedence constraint does not provide a stronger restriction than the
single-shuttle serial-workload constraint at the corresponding completion
rank, and therefore

\[
C_k^{\mathrm{mixed}}
=
C_k^{\mathrm{simple}}.
\]

Conversely, if

\[
R_k^{\mathrm{sorted}}
>
C_k^{\mathrm{simple}},
\]

the relaxation of the target-to-target in-lane precedence in the simple mode
results in a more optimistic estimate of the cumulative work. In this case,

\[
C_k^{\mathrm{mixed}}
=
R_k^{\mathrm{sorted}},
\]

and the mixed mode tightens the simple bound using the additional lane
precedence information.

Accordingly, \texttt{mixed\_shuttle\_LB} should not be interpreted as
recovering the complete whole-layout topology or as replacing the original
single-shuttle system with an actual multi-shuttle model. Instead, it retains
the single-shuttle serial-workload lower bound provided by
\texttt{simple\_shuttle\_LB} and supplements it with an independent
lane-internal precedence constraint derived from
\(R_k^{\mathrm{sorted}}\). This additional constraint can correct part of the
overly optimistic estimation introduced by the complete relaxation of
target-to-target in-lane precedence in the simple mode.

At the same time, the mixed mode still does not preserve the complete
lane-wise combinatorial consistency across different storage lanes. When
\(R_k^{\mathrm{sorted}}\) is calculated, all \(R_{\ell,r}\) values are treated
as independent lane-local bounds, and the \(R\) values associated with
different completion ranks are not required to jointly correspond to a
globally consistent selection of target counts from the individual storage
lanes. Therefore, the mixed mode preserves local lane-precedence information
while still relaxing the complete combinatorial structure of these local
precedence constraints at the whole-layout level. This remaining relaxation
constitutes the principal distinction between the mixed mode and
\texttt{topology\_shuttle\_LB}.

From the perspective of constraint relaxation,
\texttt{mixed\_shuttle\_LB} can thus be interpreted as the combination of two
complementary lower bounds:

\[
\begin{aligned}
C_k^{\mathrm{simple}}
&:
\quad
\text{single-shuttle serial-workload constraint},
\\
R_k^{\mathrm{sorted}}
&:
\quad
\text{lane-internal precedence constraint}.
\end{aligned}
\]

These two constraints are combined through

\[
C_k^{\mathrm{mixed}}
=
\max
\left\{
C_k^{\mathrm{simple}},
R_k^{\mathrm{sorted}}
\right\},
\]

such that the tighter of the two lower-bound constraints is retained for each
completion rank. The mixed mode can therefore be understood as a
completion-rank-specific tightening of the simple mode using the lane-local
precedence-constrained cumulative work lower bounds, rather than as a
replacement of the original single-shuttle lower bound by an implicitly
multi-shuttle model.

\subsubsection{Topological shuttle lower-bound mode}

The \texttt{topology\_shuttle\_LB} mode further preserves the complete
precedence-constrained cumulative work structure within each storage lane on top
of the relaxations adopted by \texttt{simple\_shuttle\_LB} and
\texttt{mixed\_shuttle\_LB}. Unlike the mixed mode, which supplements each
completion rank independently using the globally sorted
\(R_k^{\mathrm{sorted}}\), the topology mode does not detach individual
\(R_{\ell,r}\) values from the storage lanes to which they belong. Instead, the
association between each \(R_{\ell,r}\), its corresponding storage lane, and
the target-item depth order within that lane is retained throughout the
layout-level lower-bound construction. Consequently, when constructing the
cumulative work lower bound for completion rank \(k\), the selected target items
must jointly form a combination that is consistent with the precedence
relations in all involved storage lanes.

For an arbitrary storage lane \(\ell\), let
\(n_{\ell}^{\mathrm{target}}\) denote the number of currently unfinished target
items in this lane. These target items are ordered from the lane entrance toward
the deeper end as

\[
i^{\mathrm{target}}_{\ell,1},
\;
i^{\mathrm{target}}_{\ell,2},
\;
\ldots,
\;
i^{\mathrm{target}}_{\ell,n_{\ell}^{\mathrm{target}}}.
\]

The corresponding precedence-constrained cumulative work lower bounds derived
in the previous subsection are

\[
R_{\ell,1},
\;
R_{\ell,2},
\;
\ldots,
\;
R_{\ell,n_{\ell}^{\mathrm{target}}},
\]

with

\[
R_{\ell,r}
=
\sum_{h=1}^{r}
q_{\ell,h}.
\]

In addition,

\[
R_{\ell,0}=0
\]

is defined.

Therefore, selecting \(r\) completed target items from storage lane \(\ell\)
does not mean selecting any arbitrary \(r\) values of \(q_{\ell,h}\). Instead,
the selected target items must be the first \(r\) target items counted from the
lane entrance, and the corresponding minimum required work is
\(R_{\ell,r}\). In other words, \(R_{\ell,r}\) restricts the feasible
target-selection structure within storage lane \(\ell\) to a
precedence-consistent prefix structure: if a deeper target item
\(i^{\mathrm{target}}_{\ell,r}\) is included among the completed target items,
all target-related required work associated with the target items located in
front of it must also be included.

This constitutes the main distinction between the topology mode and the two
preceding modes. In the simple mode, all \(q_{\ell,r}\) values can be globally
sorted across the entire layout, such that a deeper target item may be treated
as an early completion candidate independently of the target-related work
located in front of it. In the mixed mode, although \(R_{\ell,r}\) provides an
additional lane-precedence constraint for each completion rank, the individual
\(R_{\ell,r}\) values are still globally sorted and used independently, and the
\(R\) values associated with different completion ranks are not required to
jointly correspond to a lane-wise consistent selection of target counts. The
topology mode removes this remaining relaxation. For each completion rank
\(k\), the number of target items selected from every storage lane is explicitly
considered, and the selected target items in each lane must respect the original
precedence order.

Let \(n^{\mathrm{lane}}\) denote the number of storage lanes that contain
target-related work relevant to the current lower-bound calculation. Define

\[
n_{\ell}^{\mathrm{selected}}
\]

as the number of target items selected for completion from storage lane
\(\ell\) when constructing a lower bound for a particular completion rank.
Then,

\[
n_{\ell}^{\mathrm{selected}}
\in
\left\{
0,1,\ldots,n_{\ell}^{\mathrm{target}}
\right\}.
\]

Selecting \(n_{\ell}^{\mathrm{selected}}\) target items from storage lane
\(\ell\) means completing the first
\(n_{\ell}^{\mathrm{selected}}\) target items closest to the lane entrance.
Accordingly, the minimum required work contributed by storage lane \(\ell\) is

\[
R_{\ell,n_{\ell}^{\mathrm{selected}}}.
\]

For completion rank \(k\) at the layout level, the total number of selected
target items over all storage lanes must be exactly \(k\), i.e.,

\[
\sum_{\ell=1}^{n^{\mathrm{lane}}}
n_{\ell}^{\mathrm{selected}}
=
k.
\]

For any lane-wise target-count combination satisfying this condition, the
corresponding cumulative required-work lower bound is

\[
\sum_{\ell=1}^{n^{\mathrm{lane}}}
R_{\ell,n_{\ell}^{\mathrm{selected}}}.
\]

This summation continues to reflect the single-shuttle serial workload
constraint of the actual system. Although the processing sequence among
different storage lanes is not predetermined, once the number of target items
to be completed from each lane is specified, the corresponding minimum
required work from all selected lanes must ultimately be performed by the same
shuttle. Hence, the required work contributed by different storage lanes must
be accumulated rather than treated as if it could be executed in parallel.

However, to ensure that the resulting quantity remains a lower bound
independent of any particular scheduling policy, the topology mode does not
prescribe which storage lane should be processed first, nor does it fix the
actual retrieval sequence across different storage lanes. Instead, for a given
completion rank \(k\), the most favorable lane-wise target-count combination is
selected among all combinations satisfying

\[
\sum_{\ell=1}^{n^{\mathrm{lane}}}
n_{\ell}^{\mathrm{selected}}
=
k.
\]

Accordingly, the cumulative work lower bound under
\texttt{topology\_shuttle\_LB} is defined as

\begin{equation}
C_k^{\mathrm{topology}}
=
\min_{
\substack{
0
\leq
n_{\ell}^{\mathrm{selected}}
\leq
n_{\ell}^{\mathrm{target}}
\\
\sum_{\ell=1}^{n^{\mathrm{lane}}}
n_{\ell}^{\mathrm{selected}}
=
k
}
}
\;
\sum_{\ell=1}^{n^{\mathrm{lane}}}
R_{\ell,n_{\ell}^{\mathrm{selected}}}.
\end{equation}

Here, \(C_k^{\mathrm{topology}}\) denotes the minimum cumulative amount of work
that the single shuttle must perform, from the current state, before any
\(k\) target items can be completed while preserving the necessary precedence
structure within every storage lane;
\(n^{\mathrm{lane}}\) denotes the number of storage lanes involved in the
current lower-bound calculation;
\(n_{\ell}^{\mathrm{target}}\) denotes the number of currently unfinished target
items in storage lane \(\ell\);
\(n_{\ell}^{\mathrm{selected}}\) denotes the number of target items selected for
completion from storage lane \(\ell\); and
\(R_{\ell,n_{\ell}^{\mathrm{selected}}}\) denotes the
precedence-constrained cumulative work lower bound required to complete the
first \(n_{\ell}^{\mathrm{selected}}\) target items in that lane.

For example, consider a layout containing three storage lanes and suppose that
the cumulative work lower bound for the second completion rank is to be
constructed, i.e., \(k=2\). Provided that the corresponding lanes contain
sufficient target items, the feasible lane-wise target-count combinations
include

\[
(2,0,0),
\quad
(0,2,0),
\quad
(0,0,2),
\quad
(1,1,0),
\quad
(1,0,1),
\quad
(0,1,1).
\]

If

\[
(2,0,0)
\]

is selected, both completed target items originate from the first storage lane.
Because the precedence relation within this lane must be respected, the
corresponding cumulative required work is

\[
R_{1,2}.
\]

If

\[
(1,1,0)
\]

is selected, one target item closest to the lane entrance is completed from
each of the first and second storage lanes, and the corresponding cumulative
required work is

\[
R_{1,1}
+
R_{2,1}.
\]

The remaining feasible combinations can be evaluated in the same manner.
Consequently, the topology mode selects the minimum value among all feasible
combinations:

\[
C_2^{\mathrm{topology}}
=
\min
\left\{
R_{1,2},
R_{2,2},
R_{3,2},
R_{1,1}+R_{2,1},
R_{1,1}+R_{3,1},
R_{2,1}+R_{3,1}
\right\},
\]

where only combinations permitted by the current number of target items in each
storage lane are included.

This minimization has two simultaneous but conceptually distinct effects. On
the one hand, for each individual storage lane, only one valid state from

\[
R_{\ell,0},
R_{\ell,1},
\ldots,
R_{\ell,n_{\ell}^{\mathrm{target}}}
\]

may be selected. Therefore, target-related work located closer to the lane
entrance cannot be bypassed, as in the simple mode, and individual globally
sorted \(R_k^{\mathrm{sorted}}\) values are not used independently, as in the
mixed mode. Hence, every candidate combination preserves the LIFO-type access
precedence within all selected storage lanes. On the other hand, the allocation
of the \(k\) completed target items across different storage lanes is still
chosen optimistically through global minimization. This avoids prescribing any
specific scheduling sequence across the storage lanes and thereby preserves the
optimistic nature required for a valid layout-level lower bound.

From another perspective, the topology mode searches for the minimum cumulative
required work among all precedence-consistent subsets of \(k\) target items.
In the simple mode, any \(k\) target items in the entire layout may be selected.
In the topology mode, an additional restriction is imposed: if the selected
subset contains the \(r\)-th target item
\(i^{\mathrm{target}}_{\ell,r}\) in storage lane \(\ell\), then the target items
indexed \(1,\ldots,r-1\) in the same lane must also be represented through the
corresponding preceding required work. Therefore, the set of feasible
target-item combinations under the topology mode is a subset of that under the
simple mode. Using the same \(q_{\ell,r}\) values, it follows that

\[
C_k^{\mathrm{topology}}
\geq
C_k^{\mathrm{simple}}.
\]

This relation does not imply that the topology mode overestimates the actual
cumulative work. Rather, the simple mode permits additional idealized
combinations that violate the physical in-lane precedence, whereas the topology
mode removes these infeasible combinations and thereby yields a tighter lower
bound.

The topology mode also retains the lane-precedence information exploited by the
mixed mode. For any feasible lane-wise target-count combination corresponding
to completion rank \(k\), the resulting cumulative required work cannot be
smaller than the \(R_{\ell,r}\) value associated with any completed target item
contained in that combination. Therefore, the cumulative work must also satisfy
the lower-bound constraint implied by the globally \(k\)-th smallest
precedence-constrained cumulative work lower bound,

\[
C_k^{\mathrm{topology}}
\geq
R_k^{\mathrm{sorted}}.
\]

Together with

\[
C_k^{\mathrm{topology}}
\geq
C_k^{\mathrm{simple}},
\]

this yields

\[
C_k^{\mathrm{topology}}
\geq
\max
\left\{
C_k^{\mathrm{simple}},
R_k^{\mathrm{sorted}}
\right\}.
\]

Since the mixed mode is defined as

\[
C_k^{\mathrm{mixed}}
=
\max
\left\{
C_k^{\mathrm{simple}},
R_k^{\mathrm{sorted}}
\right\},
\]

the three cumulative work lower bounds satisfy

\[
C_k^{\mathrm{simple}}
\leq
C_k^{\mathrm{mixed}}
\leq
C_k^{\mathrm{topology}}.
\]

These inequalities describe the progressive tightening of the cumulative work
lower bound as more storage-lane precedence information is preserved. They
should not be interpreted as strict inequalities, because two or all three
modes may yield identical values for particular system states and completion
ranks.

In practice, explicitly enumerating all combinations

\[
\left(
n_{1}^{\mathrm{selected}},
n_{2}^{\mathrm{selected}},
\ldots,
n_{n^{\mathrm{lane}}}^{\mathrm{selected}}
\right)
\]

becomes increasingly expensive as the number of storage lanes and target items
grows. Therefore, the current implementation computes the topology lower bound
using dynamic programming. The sequences

\[
R_{\ell,0},
R_{\ell,1},
\ldots,
R_{\ell,n_{\ell}^{\mathrm{target}}}
\]

are merged one storage lane at a time. When a new storage lane is introduced,
the algorithm considers the possible numbers of target items selected from that
lane, evaluates their contribution to each completion rank, and retains the
minimum cumulative required work for every rank. After all storage lanes have
been processed, the procedure simultaneously obtains

\[
C_1^{\mathrm{topology}},
C_2^{\mathrm{topology}},
\ldots,
C_{n^{\mathrm{target}}}^{\mathrm{topology}}.
\]

This dynamic-programming procedure improves the computational efficiency of the
lane-wise combination problem without changing the mathematical definition of
the lower bound.

It should be emphasized that the term ``topology'' in
\texttt{topology\_shuttle\_LB} does not imply that all physical routing and
traffic constraints of the actual warehouse layout are preserved. The main
additional topology information retained by this mode is the target-level
precedence structure within each storage lane induced by the current item
arrangement and the LIFO-type access constraint. The underlying
\(\tau^{\mathrm{retrieval}}\), \(\tau^{\mathrm{relocation}}\),
\(q_{\ell,r}\), and \(R_{\ell,r}\) still inherit the operation-level
relaxations introduced previously. For example, actual cross-lane
repositioning, specific relocation destinations, cross-aisle congestion, route
conflicts, waiting times, and other factors that may increase the actual
execution time are not fully represented. In addition, the actual switching
sequence and travel continuity across different storage lanes are not fixed;
instead, the most favorable lane-wise combination is retained at the
layout-level lower-bound stage.

Therefore, \texttt{topology\_shuttle\_LB} should not be interpreted as a full
simulation of the actual single-shuttle scheduling process. Rather, it preserves
the storage-lane internal precedence structure more completely than the simple
and mixed modes while retaining the relaxations required for constructing a
valid lower bound. Its main characteristics can be summarized as

\[
\begin{aligned}
&\text{retaining target-specific unavoidable work},\\
&\text{retaining single-shuttle serial workload accumulation},\\
&\text{fully preserving the precedence-consistent prefix structure within each storage lane},\\
&\text{selecting the most favorable target-count combination across different storage lanes}.
\end{aligned}
\]

Accordingly, the progressive relationship among the three layout-level modes
can be summarized as

\[
\begin{aligned}
\text{Simple:}\quad&
\text{global selection of independent } q,\\
\text{Mixed:}\quad&
\text{simple bound}
+
\text{an independent } R\text{-based precedence constraint},\\
\text{Topology:}\quad&
\text{joint lane-wise combination of complete } R\text{ structures}.
\end{aligned}
\]

Thus, \texttt{topology\_shuttle\_LB} removes the lane-wise combinatorial
relaxation introduced by globally and independently sorting the
\(R_{\ell,r}\) values in the mixed mode. By jointly optimizing over the complete
precedence-constrained cumulative work structures of all storage lanes, it
constructs, for each completion rank \(k\), the cumulative work lower bound
\(C_k^{\mathrm{topology}}\) that preserves the most complete storage-lane
topology information among the three proposed modes.

\subsection{State-level total-tardiness lower bound derivation}

In the preceding layout-level lower-bound derivation, the cumulative work lower
bounds corresponding to each completion rank \(k\) were constructed under the
three proposed lower-bound modes, namely

\[
C_k^{\mathrm{simple}},
\qquad
C_k^{\mathrm{mixed}},
\qquad
C_k^{\mathrm{topology}}.
\]

These quantities represent the minimum cumulative amount of work that is
unavoidable, under the corresponding constraint relaxation, before any \(k\)
of the current target items can be completed from the current system state.
However, \(C_k^m\) only characterizes the minimum work required to complete
different numbers of target items and does not yet explicitly incorporate the
due dates of the target items. Therefore, it cannot directly represent a lower
bound on the optimization objective considered in this study, namely total
tardiness. To derive a state-level total-tardiness lower bound, the
completion-rank cumulative work lower bounds are further combined with the due
dates of the currently unfinished target items.

Let the current system state be denoted by \(s\), and let

\[
t(s)
\]

denote the current system time. Assume that the current state contains

\[
n^{\mathrm{target}}(s)
\]

unfinished target items. The due dates of these target items are sorted in
ascending order as

\[
d_{(1)}^{\mathrm{due}}(s)
\leq
d_{(2)}^{\mathrm{due}}(s)
\leq
\cdots
\leq
d_{\left(n^{\mathrm{target}}(s)\right)}^{\mathrm{due}}(s).
\]

Here, \(d_{(k)}^{\mathrm{due}}(s)\) denotes the \(k\)-th smallest due date
among all unfinished target items in state \(s\). The parenthesized subscript
\((k)\) denotes the position of the corresponding due date in the sorted
due-date sequence, while the superscript \(\mathrm{due}\) identifies the
quantity as a due date.

For any lower-bound mode

\[
m
\in
\left\{
\mathrm{simple},
\mathrm{mixed},
\mathrm{topology}
\right\},
\]

\(C_k^m(s)\) denotes the cumulative work lower bound required to complete any
\(k\) of the currently unfinished target items from state \(s\). Consequently,
even under the most favorable subsequent scheduling policy, the \(k\)-th
target-item completion cannot occur earlier than

\[
t(s)
+
C_k^m(s).
\]

Here, \(t(s)\) represents the amount of system time that has already elapsed,
whereas \(C_k^m(s)\) represents the minimum remaining cumulative work required
from the current state before the \(k\)-th target-item completion can occur.
Their sum therefore constitutes a lower bound on the earliest possible
completion time associated with completion rank \(k\).

It should be noted that the layout-level lower-bound construction does not
predetermine which specific target item will occupy completion rank \(k\).
Accordingly, \(C_k^m(s)\) cannot be directly associated with the due date of a
predefined target item. To preserve the optimistic nature required for a valid
lower bound, the completion ranks and the due dates of the currently unfinished
target items must be paired in the most favorable manner with respect to total
tardiness.

Since the completion-time lower bounds are ordered by completion rank, an
optimistic assignment is obtained by pairing earlier completion ranks with
smaller due dates. Therefore, the \(k\)-th completion-time lower bound

\[
t(s)
+
C_k^m(s)
\]

is paired with the \(k\)-th smallest due date

\[
d_{(k)}^{\mathrm{due}}(s).
\]

The unavoidable tardiness lower bound associated with completion rank \(k\) is
then given by

\[
\max
\left\{
0,\;
t(s)
+
C_k^m(s)
-
d_{(k)}^{\mathrm{due}}(s)
\right\}.
\]

In this expression, \(t(s)+C_k^m(s)\) denotes the earliest possible completion
time of the \(k\)-th completed target item under lower-bound mode \(m\), while
\(d_{(k)}^{\mathrm{due}}(s)\) denotes the \(k\)-th smallest due date among the
currently unfinished target items. If the earliest possible completion time is
no later than the corresponding due date, no tardiness is unavoidable for that
completion rank, and the contribution is therefore truncated to zero through
the maximum operator.

In addition to the tardiness that may unavoidably be incurred by the unfinished
target items, the current state may already contain target items that have been
completed and unloaded at an Input/Output (I/O) point. For these completed
tasks, their actual tardiness has already been realized and cannot be affected
by any subsequent scheduling decision. Therefore, this part of the total
tardiness does not require further lower-bound estimation and is directly
included using its realized value.

Let

\[
T^{\mathrm{done}}(s)
\]

denote the cumulative actual tardiness of all target items that have already
been completed in state \(s\). The state-level total-tardiness lower bound under
lower-bound mode \(m\) is then defined as
\begin{equation}
LB^m(s)
=
T^{\mathrm{done}}(s)
+
\sum_{k=1}^{n^{\mathrm{target}}(s)}
\max
\left\{
0,\;
t(s)
+
C_k^m(s)
-
d_{(k)}^{\mathrm{due}}(s)
\right\}.
\end{equation}

Here, \(LB^m(s)\) denotes the total-tardiness lower bound of state \(s\) under
lower-bound mode \(m\);
\(T^{\mathrm{done}}(s)\) denotes the realized cumulative tardiness of the
already completed target items;
\(n^{\mathrm{target}}(s)\) denotes the number of currently unfinished target
items;
\(t(s)\) denotes the current system time;
\(C_k^m(s)\) denotes the cumulative work lower bound required to complete any
\(k\) of the current target items under mode \(m\); and
\(d_{(k)}^{\mathrm{due}}(s)\) denotes the \(k\)-th smallest due date among the
currently unfinished target items.

Accordingly, the state-level total-tardiness lower bound can be decomposed into
two conceptually different components:

\[
LB^m(s)
=
\underbrace{
T^{\mathrm{done}}(s)
}_{\text{realized and irreversible tardiness}}
+
\underbrace{
\sum_{k=1}^{n^{\mathrm{target}}(s)}
\max
\left\{
0,\;
t(s)
+
C_k^m(s)
-
d_{(k)}^{\mathrm{due}}(s)
\right\}
}_{\text{lower bound on future unavoidable tardiness}}.
\]

The first component represents the scheduling cost that has already been
realized and can no longer be changed, whereas the second component provides an
optimistic estimate of the minimum additional tardiness that is unavoidable
from the current state. This estimate jointly reflects the current warehouse
configuration, the current system time, the unfinished target items and their
due dates, as well as the workload and precedence information retained by the
corresponding lower-bound mode.

Therefore, \(LB^m(s)\) characterizes not only the tardiness that has already
accumulated, but also the minimum future tardiness implied by the current
system state before all remaining target items have actually been completed.
This enables the lower bound to reflect changes in the future scheduling
difficulty caused by changes in the current layout even before those changes
are fully manifested as realized tardiness.

The three proposed lower-bound modes employ the same transformation from
completion-rank cumulative work lower bounds to the state-level total-tardiness
lower bound. Their only difference at this stage lies in the corresponding
values of

\[
C_k^{\mathrm{simple}},
\qquad
C_k^{\mathrm{mixed}},
\qquad
C_k^{\mathrm{topology}}.
\]

Accordingly, three state-level total-tardiness lower bounds are obtained:

\[
LB^{\mathrm{simple}}(s),
\qquad
LB^{\mathrm{mixed}}(s),
\qquad
LB^{\mathrm{topology}}(s).
\]

As established in the preceding derivation,

\[
C_k^{\mathrm{simple}}(s)
\leq
C_k^{\mathrm{mixed}}(s)
\leq
C_k^{\mathrm{topology}}(s).
\]

Since

\[
\max
\left\{
0,\;
t+C-d
\right\}
\]

is monotonically nondecreasing with respect to \(C\), the corresponding
state-level total-tardiness lower bounds satisfy

\[
LB^{\mathrm{simple}}(s)
\leq
LB^{\mathrm{mixed}}(s)
\leq
LB^{\mathrm{topology}}(s).
\]

These inequalities indicate that the state-level total-tardiness lower bound
can be progressively tightened as more storage-lane precedence information is
preserved. They should not be interpreted as strict inequalities, since two or
all three lower-bound modes may yield identical values for particular system
states when the additional precedence constraints do not further increase the
unavoidable tardiness estimate.

\subsection{Reward construction based on the state-level lower bound}

In the preceding subsection, the state-level total-tardiness lower bounds under
the three proposed lower-bound modes were constructed as

\[
LB^{\mathrm{simple}}(s),
\qquad
LB^{\mathrm{mixed}}(s),
\qquad
LB^{\mathrm{topology}}(s).
\]

These state-level lower bounds incorporate not only the realized tardiness that
has already occurred and can no longer be changed, but also an estimate of the
future unavoidable tardiness based on the due dates of the currently unfinished
target items and the corresponding completion-rank cumulative work lower bounds
\(C_k^m(s)\). Therefore, compared with a reward constructed solely from realized
tardiness, \(LB^m(s)\) can reflect the potential influence of the current
scheduling decision on future total tardiness at an earlier stage.

Specifically, in a multi-deep storage system, an upper-level decision may first
change the accessibility of target items, the blocking relationships, the
necessary relocation work, or the precedence structure within a storage lane,
without immediately completing a target item. Consequently, if the reward is
constructed only from changes in realized tardiness, a retrieval or relocation
decision that substantially improves the future scheduling state may not receive
timely feedback at the current decision point. In contrast, because
\(LB^m(s)\) changes with the current layout, the shuttle carrying state, the
unfinished target items, and the remaining required work, the difference
between the lower bounds of consecutive upper-level decision states can be used
to convert the effect of the current action on future unavoidable total
tardiness into more immediate feedback.

To provide a unified representation of the three lower-bound modes, let

\[
m
\in
\left\{
\mathrm{simple},
\mathrm{mixed},
\mathrm{topology}
\right\}.
\]

Let \(s_t\) denote the system state at the \(t\)-th upper-level decision point,
and let the upper-level agent select action \(a_t\) in this state. After the
retrieval or relocation subtask associated with this upper-level action has been
completed, the system reaches the next upper-level decision state
\(s_{t+1}\). Without considering the counterfactual correction associated with
dynamic target activation at this stage, the reward for the upper-level action
is constructed from the change in the state-level total-tardiness lower bound
between these two states.

Since \(LB^m(s)\) represents a cost-like quantity to be minimized, a smaller
lower bound corresponds to a more favorable system state. To make the direction
of state evaluation consistent with the reward-maximization objective in
reinforcement learning, the potential function of state \(s\) under
lower-bound mode \(m\) is defined as

\[
\Phi^m(s)
=
-LB^m(s).
\]

Here, \(\Phi^m(s)\) denotes the potential value of state \(s\) under mode \(m\),
and \(LB^m(s)\) denotes the corresponding state-level total-tardiness lower
bound. The negative sign reverses the direction of the cost quantity such that
a lower total-tardiness lower bound corresponds to a higher state potential.
Accordingly, for any two states \(s^A\) and \(s^B\), if

\[
LB^m(s^A)
<
LB^m(s^B),
\]

then

\[
\Phi^m(s^A)
>
\Phi^m(s^B).
\]

Thus, a state with a lower amount of unavoidable total tardiness is assigned a
higher potential value.

In general, the potential difference between two consecutive states can be
expressed as

\[
\gamma
\Phi^m(s_{t+1})
-
\Phi^m(s_t),
\]

where \(\gamma\) denotes the discount factor. In the current training setting,

\[
\gamma=1.
\]

Therefore, the potential difference between two consecutive upper-level
decision states reduces to

\[
\Phi^m(s_{t+1})
-
\Phi^m(s_t).
\]

Substituting

\[
\Phi^m(s)
=
-LB^m(s)
\]

gives

\[
\Phi^m(s_{t+1})
-
\Phi^m(s_t)
=
LB^m(s_t)
-
LB^m(s_{t+1}).
\]

Accordingly, the reward assigned to the \(t\)-th upper-level action under
lower-bound mode \(m\) is defined as

\[
r_t^m
=
LB^m(s_t)
-
LB^m(s_{t+1}).
\]

Here, \(r_t^m\) denotes the lower-bound-based reward assigned to the
\(t\)-th upper-level action; \(LB^m(s_t)\) denotes the state-level
total-tardiness lower bound before the action is executed; and
\(LB^m(s_{t+1})\) denotes the state-level total-tardiness lower bound at the
next upper-level decision state after the corresponding subtask has been
completed.

If the current upper-level action reduces the state-level lower bound, i.e.,

\[
LB^m(s_{t+1})
<
LB^m(s_t),
\]

then

\[
r_t^m>0.
\]

This indicates that the action reduces the total tardiness that remains
unavoidable from the new state and therefore receives a positive reward. Such a
reduction may result from, for example, decreasing the necessary future
relocation work, improving the accessibility of target items, completing part
of the unavoidable required work, or reducing the minimum cumulative workload
associated with subsequent completion ranks.

Conversely, if

\[
LB^m(s_{t+1})
>
LB^m(s_t),
\]

then

\[
r_t^m<0,
\]

indicating that the action moves the system to a state with a larger amount of
unavoidable total tardiness. If

\[
LB^m(s_{t+1})
=
LB^m(s_t),
\]

then

\[
r_t^m=0,
\]

meaning that although the action changes the system state, it does not alter
the future unavoidable total tardiness captured by the corresponding
lower-bound mode.

Therefore, the reward does not directly evaluate whether an individual
retrieval or relocation operation immediately completes a target item. Instead,
it evaluates whether the upper-level action improves the potential of the
overall system state. Using

\[
\Phi^m(s)
=
-LB^m(s),
\]

the reward can equivalently be expressed as

\[
r_t^m
=
\Phi^m(s_{t+1})
-
\Phi^m(s_t).
\]

Hence, under the current setting of \(\gamma=1\), the lower-bound-based reward
has the mathematical form of a potential-difference reward between consecutive
upper-level decision states. It should be noted that this potential difference
is directly used as the upper-level training reward rather than being added to
the original tardiness reward as an additional shaping term. Therefore, the
proposed design is more precisely characterized as a potential-difference
reward based on the state-level total-tardiness lower bound, rather than as an
additional potential-based shaping term applied to the original reward.

The proposed reward construction also has an important cumulative property.
Consider a normally completed episode with the sequence of upper-level decision
states

\[
s_0,
s_1,
\ldots,
s_T.
\]

Without considering additional truncation penalties or counterfactual
corrections, the cumulative lower-bound-based reward over the episode is

\[
\sum_{t=0}^{T-1}
r_t^m
=
\sum_{t=0}^{T-1}
\left[
LB^m(s_t)
-
LB^m(s_{t+1})
\right].
\]

Because the intermediate lower-bound terms cancel telescopically, this
expression reduces to

\[
\sum_{t=0}^{T-1}
r_t^m
=
LB^m(s_0)
-
LB^m(s_T).
\]

This shows that the intermediate variations in the lower bound determine how
the reward is distributed among the upper-level decisions, while the total
episode reward remains equal to the difference between the initial and terminal
state lower bounds.

When an episode terminates normally, all target items have been completed, and
therefore

\[
n^{\mathrm{target}}(s_T)=0.
\]

At this point, the future unavoidable tardiness component of the state-level
lower bound disappears, such that

\[
LB^m(s_T)
=
T^{\mathrm{done}}(s_T).
\]

Since all target items have been completed at the terminal state,
\(T^{\mathrm{done}}(s_T)\) is equal to the final total tardiness of the episode,
denoted by

\[
T^{\mathrm{total}}.
\]

Hence,

\[
LB^m(s_T)
=
T^{\mathrm{total}},
\]

and therefore

\[
\sum_{t=0}^{T-1}
r_t^m
=
LB^m(s_0)
-
T^{\mathrm{total}}.
\]

For a fixed problem instance, the initial state \(s_0\) and the corresponding
lower bound \(LB^m(s_0)\) remain unchanged across different scheduling
policies. Therefore, under normal episode termination and in the absence of
additional reward terms, maximizing the cumulative lower-bound-based reward is
equivalent to minimizing the final total tardiness. At the same time, the
changes in the intermediate state lower bounds redistribute the originally
delayed total-tardiness objective into immediate feedback at multiple
upper-level decision points, thereby alleviating the delayed-reward problem in
long-horizon scheduling.

All three lower-bound modes employ the same reward construction mechanism. Their
difference lies only in the amount of precedence information retained by the
state evaluation function \(LB^m(s)\). Accordingly, the rewards are given by

\[
r_t^{\mathrm{simple}}
=
LB^{\mathrm{simple}}(s_t)
-
LB^{\mathrm{simple}}(s_{t+1}),
\]

\[
r_t^{\mathrm{mixed}}
=
LB^{\mathrm{mixed}}(s_t)
-
LB^{\mathrm{mixed}}(s_{t+1}),
\]

and

\[
r_t^{\mathrm{topology}}
=
LB^{\mathrm{topology}}(s_t)
-
LB^{\mathrm{topology}}(s_{t+1}).
\]

The corresponding potential functions are

\[
\Phi^{\mathrm{simple}}(s)
=
-LB^{\mathrm{simple}}(s),
\]

\[
\Phi^{\mathrm{mixed}}(s)
=
-LB^{\mathrm{mixed}}(s),
\]

and

\[
\Phi^{\mathrm{topology}}(s)
=
-LB^{\mathrm{topology}}(s).
\]

Consequently, the complete hierarchical construction of the lower-bound-based
reward can be summarized as

\[
\left\{
\tau^{\mathrm{retrieval}},
\tau^{\mathrm{relocation}}
\right\}
\longrightarrow
q_{\ell,r}
\longrightarrow
R_{\ell,r}
\longrightarrow
C_k^m
\longrightarrow
LB^m(s)
\longrightarrow
\Phi^m(s)
\longrightarrow
r_t^m.
\]

The lower bounds on individual retrieval and relocation operation times provide
the fundamental required-work estimates; \(q_{\ell,r}\) aggregates the work
associated with a target item and its necessary blockers; \(R_{\ell,r}\)
further incorporates the internal precedence of the storage lane;
\(C_k^m\) aggregates these local lower bounds into layout-level cumulative work
lower bounds for different completion ranks; \(LB^m(s)\) combines them with the
current time and due dates to obtain the state-level total-tardiness lower bound;
and finally, the improvement in the state potential is converted into the
immediate reward assigned to each upper-level action.
\subsection{Counterfactual reward shaping}

To reduce the random fluctuations introduced into upper-level action rewards by unpredictable dynamic target activation, this work further adopts counterfactual state evaluation for the three lower-bound-based reward modes and uses it to shape the basic state-level lower-bound reward. In this work, \emph{counterfactual} specifically means that, at the state reached after the current upper-level transition, the actual physical warehouse layout, current time, shuttle state, and already realized tardiness are preserved, while the items newly activated during the current transition are temporarily treated as non-target items only for the lower-bound evaluation used to calculate the current reward.

This design addresses the specific reward-assignment problem introduced by dynamic order-driven target activation. From the perspective of the current upper-level decision point, target activations occurring during the execution of an upper-level action constitute unpredictable exogenous state changes caused by the dynamic order process. An upper-level retrieval or relocation action generally requires multiple discrete time steps to complete. Therefore, between the action-start state $s_t$ and the subsequent upper-level decision state $s_{t+1}$, one or more non-target items may be dynamically activated and assigned due dates. If the lower bound of the actual state $s_{t+1}$ is directly used to evaluate the current action, the resulting reward reflects not only the effect of the action on the warehouse layout, blocking relationships, and remaining required work, but also the change in the target-item set caused by newly arriving orders. Consequently, even identical or similar scheduling actions may receive substantially different rewards depending on the target activations that happen to occur during their execution.

Consistent with the state-level total-tardiness lower bound derived above, for any lower-bound mode

\[
m \in
\{\mathrm{simple},\mathrm{mixed},\mathrm{topology}\},
\]

the lower bound of a state $s$ is defined as

\[
LB^m(s)
=
T^{\mathrm{done}}(s)
+
\sum_{k=1}^{n^{\mathrm{target}}(s)}
\max
\left\{
0,\,
t(s)+C_k^m(s)-d_{(k)}^{\mathrm{due}}(s)
\right\}.
\]

Here, $LB^m(s)$ denotes the state-level total-tardiness lower bound of state $s$ under lower-bound mode $m$; $T^{\mathrm{done}}(s)$ denotes the realized cumulative tardiness of the already completed target items; $n^{\mathrm{target}}(s)$ denotes the number of currently unfinished target items; $t(s)$ denotes the current system time; $C_k^m(s)$ denotes the cumulative work lower bound required to complete any $k$ of the currently unfinished target items under mode $m$; and $d_{(k)}^{\mathrm{due}}(s)$ denotes the $k$th smallest due date among the currently unfinished target items.

Without counterfactual reward shaping, after the upper-level agent selects action $a_t$ in state $s_t$ and the corresponding retrieval or relocation subtask is completed, the system reaches the next upper-level decision state $s_{t+1}$. The basic lower-bound-based reward is

\[
r_t^m
=
LB^m(s_t)-LB^m(s_{t+1}),
\]

where $r_t^m$ denotes the reward assigned to the $t$th upper-level action under lower-bound mode $m$. This reward measures the change in the state-level total-tardiness lower bound between the states before and after the execution of the current upper-level action.

Let

\[
\Delta\mathcal{T}_t
\]

denote the set of items that are newly activated as target items during the upper-level transition from $s_t$ to $s_{t+1}$, that is, from the beginning of the current upper-level decision until the associated retrieval or relocation subtask is completed.

Under the counterfactual mode, the actual physical state reached after the transition is retained, but a counterfactual state evaluation is performed solely for the current reward calculation. Specifically, for each item in $\Delta\mathcal{T}_t$, its actual physical location and its influence on the warehouse topology and blocking relationships are preserved, while its newly assigned target status is temporarily ignored during the lower-bound calculation. The corresponding counterfactual state-level lower bound is denoted by

\[
LB_{\mathrm{cf}}^m
\left(
s_{t+1}\mid\Delta\mathcal{T}_t
\right).
\]

Here, the subscript $\mathrm{cf}$ denotes the counterfactual state evaluation. This quantity does not correspond to a separate physical warehouse state generated through an alternative environment transition. Instead, it represents an alternative lower-bound evaluation of the actual state $s_{t+1}$ in which the target activations that occurred during the current transition are temporarily excluded from the target-item set.

Importantly, the items whose newly assigned target status is temporarily ignored are not removed from the warehouse. Therefore, if an item in $\Delta\mathcal{T}_t$ remains in a storage lane and blocks another active target item at state $s_{t+1}$, it is still treated as a non-target blocker in the counterfactual state evaluation and contributes to the corresponding relocation required-work lower bound. Hence, the counterfactual state evaluation removes only the immediate influence of the newly assigned target-task attributes on the reward of the current transition, while preserving the physical position and blocking effect of these items.

Accordingly, the counterfactual reward is defined as

\[
r_{t,\mathrm{cf}}^m
=
LB^m(s_t)
-
LB_{\mathrm{cf}}^m
\left(
s_{t+1}\mid\Delta\mathcal{T}_t
\right).
\]

Here, $r_{t,\mathrm{cf}}^m$ denotes the reward assigned to the $t$th upper-level action under lower-bound mode $m$ after applying counterfactual reward shaping. The starting point of the reward remains the lower bound of the actual pre-action state $s_t$, whereas the post-action term is replaced by the lower bound obtained through the counterfactual state evaluation of $s_{t+1}$.

Equivalently, the counterfactual reward can be written as

\[
r_{t,\mathrm{cf}}^m
=
r_t^m
+
F_{t,\mathrm{cf}}^m,
\]

where the counterfactual shaping term is

\[
F_{t,\mathrm{cf}}^m
=
LB^m(s_{t+1})
-
LB_{\mathrm{cf}}^m
\left(
s_{t+1}\mid\Delta\mathcal{T}_t
\right).
\]

Here, $F_{t,\mathrm{cf}}^m$ denotes the reward correction associated with target activations occurring during the current upper-level transition. This term separates the difference between the actual state evaluation and the counterfactual state evaluation from the direct evaluation of the current action, so that the resulting reward more specifically reflects the scheduling effect of the upper-level action itself.

If no new target activation occurs during the current upper-level transition,

\[
\Delta\mathcal{T}_t=\varnothing,
\]

then the counterfactual state evaluation is identical to the actual state evaluation, i.e.,

\[
LB_{\mathrm{cf}}^m
\left(
s_{t+1}\mid\varnothing
\right)
=
LB^m(s_{t+1}),
\]

and therefore

\[
F_{t,\mathrm{cf}}^m=0
\]

and

\[
r_{t,\mathrm{cf}}^m=r_t^m.
\]

Thus, the counterfactual reward shaping modifies the immediate reward only when dynamic target activation actually occurs during the execution of an upper-level action.

The rationale behind this design is to distinguish, within the reward evaluation of the current transition, between state changes that can be influenced by the current scheduling action and unpredictable exogenous changes caused by random order arrivals. The counterfactual state evaluation does not remove the elapsed time consumed by the current action, nor does it remove the actual effects of the action on the shuttle position, item positions, blocking relationships, or remaining required work. Therefore, an inefficient action that consumes excessive time, increases future relocation work, or creates a less favorable blocking structure is still evaluated accordingly through the lower bound. The only component temporarily excluded is the immediate effect of newly introduced target-task attributes arising from target activations during the current transition.

Moreover, the counterfactual state evaluation is used only for the reward calculation of the current transition. The actual next state remains $s_{t+1}$, in which the items in $\Delta\mathcal{T}_t$ retain their actual target status and assigned due dates. These newly activated target items therefore participate normally in the state representation, action space, and lower-bound calculation of subsequent upper-level decisions. In particular, the starting lower bound of the next transition is again evaluated using the actual state,

\[
LB^m(s_{t+1}),
\]

rather than

\[
LB_{\mathrm{cf}}^m
\left(
s_{t+1}\mid\Delta\mathcal{T}_t
\right).
\]

Therefore, the proposed reward shaping does not modify the actual evolution of the dynamic target-activation process and does not cause newly arriving orders to be ignored in subsequent decisions. Its role is limited to preventing the lower-bound change caused by unpredictable target activations during the execution of an upper-level action from being directly attributed to that action, thereby reducing the random fluctuations introduced by the dynamic order process into the action-level reward signal.
\section{Der Vortrag}\label{sec:Vortrag}
Die Folien können in \LaTeX, Powerpoint, CorelDRAW oder von
Hand erstellt werden. Eine Vorlage in Powerpoint ist unter \verb"_Vorlagen_INM" abgelegt. Folgende Punkte sind generell zu beachten:
\begin{itemize}
 \item Im Vorfeld des Vortrages ist mit dem Hauptbetreuer abzuklären wie viel Sprechzeit erlaubt ist und wann der Vortrag stattfindet. Die Sprechzeit darf auf keinen Fall überzogen werden.

 \item Machen Sie nicht zu viele Folien. Der Hauptbetreuer soll die Anzahl und die Qualität der Folien überprüfen.

 \item Die Titelfolie beinhaltet mindestens:
  \begin{itemize}
   \item Titel der Studienarbeit
   \item Name des Studierenden
   \item Typ der Arbeit (z.B.~Masterarbeit)
   \item Datum
   \item Universitäts- und Institutslogo
   \item Die Namen aller Betreuer
  \end{itemize}

  \item Achten Sie darauf, dass Ihre Folien (inklusive Grafiken) durch den Beamer gut angezeigt werden. Bearbeiten Sie gegebenenfalls Ihre Folien.
\end{itemize}
